% eksperymentalne tłumaczenie na włoski

%

\section{Okrąg Spiekera} % lang-pl
\section{Cerchio di Spieker} % lang-it

\begin{proposition}
[punkt i okrąg Spiekera] % lang-pl
[punto e cerchio di Spieker] % lang-it
\label{punkt_spiekera}%
    Niech $\triangle ABC$ będzie trójkątem. % lang-pl
    Środek okręgu wpisanego w trójkąt, którego wierzchołkami są środki boków trójkąta $ABC$, nazywamy jego punktem Spiekera. % lang-pl
    Sia $\triangle ABC$ un triangolo. % lang-it
    Il centro della circonferenza inscritta nel triangolo avente per vertici i punti medi dei lati del triangolo $ABC$ è detto punto di Spieker. % lang-it
\end{proposition}

O punkcie Spiekera pisze też Zetel \cite[s. 23]{zetel_2020}, chociaż nie nazywa go tak, być może dlatego, że Theodor Spieker będzie dziewiętnastowiecznym geometrą z Niemiec (Poczdamu)? % lang-pl
Anche Zetel \cite[p.~23]{zetel_2020} scrive del punto di Spieker, pur senza chiamarlo così, forse perché Theodor Spieker fu un geometra tedesco del XIX secolo (di Potsdam)? % lang-it
\index[persons]{Spieker, Theodor}

https://www.deltami.edu.pl/2011/11/srodek-ciezkosci/

%